\documentclass[11pt]{article}
\usepackage[margin=1in]{geometry}
\usepackage{amsmath,amssymb,amsthm,booktabs,array}
\usepackage[colorlinks=true,linkcolor=blue,urlcolor=blue]{hyperref}
\usepackage{microtype}
\newtheorem{theorem}{Theorem}[section]
\newtheorem{lemma}[theorem]{Lemma}
\newtheorem{corollary}[theorem]{Corollary}
\newcommand{\E}{\mathbb E}
\newcommand{\Var}{\operatorname{Var}}
\newcommand{\Inf}{\operatorname{Inf}}
\newcommand{\R}{\mathbb R}
\setlength{\emergencystretch}{2em}
\title{Degree-three polynomial influence for Gibbs measures:\\
expected derivatives bypass the centered-Gram obstruction}
\author{Research continuation in the AAconjecture workspace}
\date{14 November 2026}
\begin{document}
\maketitle

\noindent\textbf{Status.} A self-contained argument below proves a
dimension-independent $c\,v^4$ lower bound on physical flip influence
for globally bounded real polynomials of degree at most three. The
Hamiltonians have arbitrary finite fields, arbitrary interaction
arities, and unrestricted coordinate support. The hypotheses are a
fixed TV Dobrushin margin and a fixed bound on incident interaction
strength; in particular, every fixed weighted-strength regime
$\omega\le\omega_0<1$ is covered. A final corollary, using the precisely
stated scalar Boolean Bohnenblust--Hille theorem, extends the method
to each fixed degree. This answers the handoff's weaker
$n$-independent-$S(v,d)$ formulation, but does not give a bound
polynomial jointly in $d$ and $1/v$, or a sharp inverse-temperature
threshold. Finite computations supplement the proofs. No independent
audit, novelty claim, or formal verification is asserted.

The prompt and handoff were read from \texttt{prompt/}; the requested
\texttt{.ultramath/} archive is absent. Its unavailable manuscripts and
scripts are not inputs to this proof. The earlier quadratic report is
preserved. The corrected Gibbs-history PDF is used to keep query
simulation distinct from physical polynomial influence; its simulation
theorem is not used here.

\section{Precise statement and conventions}

For arbitrary $n\ge1$, let
\[
 H(x)=-\sum_{|A|\ge2}J_A\prod_{i\in A}x_i-\sum_i h_ix_i,
 \qquad \mu(x)=Z^{-1}e^{-\beta H(x)},\qquad x\in\{-1,1\}^n,
\]
with finite real coefficients and $\beta\ge0$. Set
\[
 K=\max_i\beta\sum_{A\ni i,\ |A|\ge2}|J_A|,
 \qquad \omega=\max_i\beta\sum_{A\ni i,\ |A|\ge2}(|A|-1)|J_A|.
\]
Here $K$ bounds interactions, not fields. No bound is imposed on
$\max_i\beta|h_i|$, the arities, or the support of the polynomial.
Put
\[
 \theta_i=\mu(X_i=1\mid X_{-i}),\quad
 v_i=4\theta_i(1-\theta_i),\quad
 m_i=\E X_i,\quad \sigma_i^2=1-m_i^2.
\]
Full support implies $\sigma_i>0$. Define $c_{ii}=0$ and
\[
 c_{ij}=\sup_{x,z:\,x_{-j}=z_{-j}}
 |\theta_i(x_{-i})-\theta_i(z_{-i})|,
 \qquad r=\max_i\sum_j c_{ij}.
\]
This is the total-variation/probability convention, not the
conditional-mean convention, which has an extra factor two.
For the multilinear derivative
$\partial_i p=(p(+1,x_{-i})-p(-1,x_{-i}))/2$, use exactly
\[
 \Inf_i^{\rm fl}=4\E(\partial_i p)^2,\qquad
 \Inf_i^{\rm pv}=\E[v_i(\partial_i p)^2],\qquad
 \Inf_i^{\rm rs2}=2\Inf_i^{\rm pv}.
\]
Thus $\Inf_i^{\rm fl}\ge4\Inf_i^{\rm pv}=2\Inf_i^{\rm rs2}$.
All derivatives and expectations below use the same original law
$\mu$; differentiation never differentiates the measure.

\begin{theorem}[Cubic influence bound]\label{thm:cubic}
Fix $K_0<\infty$, $0\le r_0<1$, and $0<v\le1$ before choosing
$n$. Let $\mu$ be any law above with $K\le K_0$, $r\le r_0$.
Define the universal numerical constant
\[
 B=1024\,26^{3/2}.
\]
Every real multilinear $p$ with $\deg p\le3$,
$\|p\|_{\infty,\{-1,1\}^n}\le1$, and $\Var_\mu p\ge v$ has
a coordinate $i$ satisfying simultaneously
\begin{align}
 \Inf_i^{\rm fl}(p)&\ge
 \frac{e^{-40K_0}(1-r_0)^{12}}{324B^4}\,v^4,
 \label{eq:mainfl}\\
 \Inf_i^{\rm pv}(p)&\ge
 \frac{e^{-44K_0}(1-r_0)^{12}}{1296B^4}\,v^4,
 \label{eq:mainpv}\\
 \Inf_i^{\rm rs2}(p)&\ge
 \frac{e^{-44K_0}(1-r_0)^{12}}{648B^4}\,v^4.
 \label{eq:mainrs}
\end{align}
The constants do not depend on $n$, field strengths, or support size.
They are deliberately conservative, not claimed optimal.
\end{theorem}

\begin{corollary}[Weighted strength and pinnings]\label{cor:omega}
For every fixed $0\le\omega_0<1$, the theorem holds uniformly when
$\omega\le\omega_0$, with $K_0=r_0=\omega_0$. In particular,
\[
 \boxed{\omega\le\tfrac18\quad\Longrightarrow\quad
 \max_i\Inf_i^{\rm fl}(p)\ge
 \frac{e^{-5}(7/8)^{12}}{324(1024\,26^{3/2})^4}\,v^4.}
\]
The same bound holds after any pinning, provided the restricted
polynomial has variance at least $v$ under the pinned law.
\end{corollary}

\begin{corollary}[Pairwise coupling-based Dobrushin region]\label{cor:pairwise}
For pairwise Ising Hamiltonians with arbitrary signs and finite fields,
assume $\max_i\sum_j\tanh(\beta|J_{ij}|)\le1-\kappa$ for fixed
$0<\kappa\le1$. Then
\[
 \max_i\Inf_i^{\rm fl}(p)\ge
 \frac{\kappa^{32}}{324B^4(2-\kappa)^{20}}\,v^4.
\]
The simultaneous posterior-variance bound is
$\kappa^{34}v^4/[1296B^4(2-\kappa)^{22}]$.
Both statements are uniform under pinnings.
\end{corollary}

\section{The genuine roadblock in the old route}

Ordinary centered Taylor coefficients do not give a uniform upper
Gram bound for cubics. To see this with exact constants, take an
independent uniform spin $X_0$ and $m$ independent Ising pairs,
\[
 H=-a\sum_{j=1}^m X_{2j-1}X_{2j},\quad
 a=\tfrac12\log(5/4)<\tfrac18,\quad
 Y_j=X_{2j-1}X_{2j},\quad \rho=\tanh a=\tfrac19.
\]
Here $K=\omega=a$, $r=\rho$, all fields vanish, and all one-site
means vanish. For the globally bounded cubic
\[
 p_m=X_0\frac1m\sum_{j=1}^mY_j
\]
one has
\[
 \E p_m=0,\quad
 \Var p_m=\rho^2+\frac{1-\rho^2}{m},\quad
 \sum_{S\ne\emptyset}\widehat p_m(S)^2=\frac1m.
\]
The ratio is $m\rho^2+1-\rho^2\to\infty$. Centering by marginal
means changes nothing. Consequently, extending the quadratic proof
by asserting a uniform Gram bound for those centered coefficients
would be invalid, even with variance bounded below.

This is not an influence counterexample:
\[
 \Inf_0^{\rm fl}(p_m)=4\left(\rho^2+\frac{1-\rho^2}{m}\right),
 \qquad \Inf_{2j-1}^{\rm fl}(p_m)=\Inf_{2j}^{\rm fl}(p_m)=\frac4{m^2}.
\]
The coordinate $0$ already has large influence. The new coefficients
detect it through $\E\partial_0p_m=\rho$, which ordinary centered
linear coefficients miss. Thus the obstruction is bypassed below,
not assumed away.

\section{Two local estimates, proved without field restrictions}

\begin{lemma}[Dobrushin Poincare estimate]\label{lem:poincare}
If $r<1$ and $C=(1-r)^{-1}$, then every function on the cube satisfies
\[
 \Var_\mu f\le C\sum_i\E\Var(f\mid X_{-i})
 =C\sum_i\E[v_i(\partial_i f)^2].
\]
In particular, for affine $f=c+\sum_i t_ix_i$,
\[
 \Var_\mu f\le C\sum_i\sigma_i^2t_i^2.
\]
\end{lemma}
\begin{proof}
Write $P_i f=\E[f\mid X_{-i}]$ and $P=n^{-1}\sum_iP_i$.
These are operators on $L^2(\mu)$; $P_i$ are orthogonal projections,
so $P$ is self-adjoint and positive semidefinite. For
$\delta_i(f)=\max_x|f(x)-f(x^{(i)})|$, coupling the two conditional
laws of spin $j$ at inputs differing at $i$ gives
\[
 \delta_i(P_j f)\le\delta_i(f)+c_{ji}\delta_j(f)\quad(i\ne j),
 \qquad \delta_i(P_i f)=0.
\]
Indeed the coupled configurations differ at $i$, and additionally
at $j$ with probability at most $c_{ji}$. Summing yields
\[
 \sum_i\delta_i(Pf)\le
 \left(1-\frac{1-r}{n}\right)\sum_i\delta_i(f).
\]
Every nonconstant eigenfunction has positive oscillation seminorm;
its eigenvalue is at most $1-(1-r)/n$. The Dirichlet identity
$\langle f,(I-P)f\rangle=n^{-1}\sum_i\E\Var(f\mid X_{-i})$
and finite-dimensional diagonalization give the claimed inequality.
For affine $f$, use $\E v_i\le\Var X_i=\sigma_i^2$, by total variance.
\end{proof}

\begin{lemma}[Local variance comparison]\label{lem:local}
With $D=e^{4K}$, for every boundary configuration,
\[
 D^{-1}\sigma_i^2\le v_i(x_{-i})\le D\sigma_i^2.
 \tag{LC}
\]
Consequently, for every multilinear $f$ and every $i$,
\begin{align}
 \Var f&\le CD\sum_i\sigma_i^2\E(\partial_i f)^2,
 \label{eq:upper}\\
 \E f^2&\ge D^{-1}\sigma_i^2\E(\partial_i f)^2.
 \label{eq:lower}
\end{align}
\end{lemma}
\begin{proof}
The conditional log odds
$\ell_i=2\beta h_i+2\beta\sum_{A\ni i}J_A\chi_{A\setminus i}$
have oscillation at most $4K$. The marginal probability
$\bar\theta_i=\E\theta_i$ lies between the extreme conditional
probabilities, so its log odds lie in the same interval. For
$g(\ell)=4e^\ell/(1+e^\ell)^2$, one has $|(\log g)'|\le1$.
Since $g(\ell_i)=v_i$ and $g(\operatorname{logit}\bar\theta_i)
=\sigma_i^2$, (LC) follows. The upper estimate follows from
Lemma~\ref{lem:poincare}; the lower follows from
$\E f^2\ge\E\Var(f\mid X_{-i})$ and (LC).
\end{proof}

For the Gibbs specification, flipping $x_j$ changes the local field
at $i$ by at most $2\beta\sum_{A\ni i,j}|J_A|$. The derivative of
$(1+\tanh u)/2$ is at most $1/2$, so
\[
 c_{ij}\le\beta\sum_{A\ni i,j}|J_A|,\qquad r\le\omega,
 \qquad K\le\omega.
 \tag{DS}
\]
No independence, global density ratio, or uniform conditional atom
floor is asserted or needed.

\section{Expected derivatives and their generating polynomial}

For a subset $S$, let $\partial_S=\prod_{i\in S}\partial_i$,
$\sigma_S=\prod_{i\in S}\sigma_i$, and define
\[
 a_S=\sigma_S\E_\mu\partial_Sp,\qquad
 a_\emptyset=\E_\mu p.
\]
For degree two these coincide with the weighted centered coefficients
used in the earlier report. For degree three, the linear coefficient
includes covariances: generally
$\E\partial_i p\ne(\partial_i p)(m)$.

\begin{lemma}[Upper energy estimate for cubics]\label{lem:energy}
If $\deg p\le3$, then
\begin{align}
 \Var_\mu p
 &\le CD\sum_{|S|=1}a_S^2
       +2C^2D^2\sum_{|S|=2}a_S^2
       +6C^3D^2\sum_{|S|=3}a_S^2\nonumber\\
 &\le L\sum_{1\le|S|\le3}a_S^2,
 \qquad L=6C^3D^2.\label{eq:energy}
\end{align}
\end{lemma}
\begin{proof}
Apply \eqref{eq:upper} to $p$, and write
$\E(\partial_i p)^2=(\E\partial_i p)^2+\Var(\partial_i p)$.
Apply \eqref{eq:upper} again to each derivative. The second derivative
$\partial_{ij}p$ is affine, so the last assertion of
Lemma~\ref{lem:poincare} gives
\[
 \E(\partial_{ij}p)^2\le(\E\partial_{ij}p)^2+
 C\sum_{k\notin\{i,j\}}\sigma_k^2(\partial_{ijk}p)^2.
\]
The third derivatives are constants. Substitution proves the first
line: each two-element subset occurs in two orders, and each triple
in six. Repeated indices contribute zero by multilinearity.
Finally $C,D\ge1$, so every coefficient is at most $L$.
\end{proof}

\begin{lemma}[A bounded generating polynomial]\label{lem:generator}
If $\deg p\le3$ and $\|p\|_\infty\le1$, the polynomial
\[
 g(z)=\E_\mu[p(X_1+\sigma_1z_1,\ldots,X_n+\sigma_nz_n)]
      =\sum_{|S|\le3}a_Sz^S
\]
is multilinear and satisfies $\|g\|_{\infty,[-1,1]^n}\le26$.
\end{lemma}
\begin{proof}
The exact multilinear Taylor expansion at each $X$ proves the
coefficient identity; expectation is under the fixed original law.
Multilinearity also extends $|p|\le1$ to the whole cube $[-1,1]^n$.
For $y\in[-2,2]^n$, set $h(t)=p(ty/2)$. This has degree at most
three and $|h(t)|\le1$ on $[-1,1]$. Lagrange interpolation at
$-1,-1/2,1/2,1$ gives
\[
 h(2)=-\tfrac52h(-1)+6h(-1/2)-10h(1/2)+\tfrac{15}2h(1).
\]
The absolute weights sum to $26$, so $|p(y)|\le26$.
Since $X+\sigma z\in[-2,2]^n$, averaging proves the bound.
\end{proof}

\section{A self-contained cubic coefficient inequality}

\begin{lemma}[Elementary scalar coefficient estimate]\label{lem:bh}
For every real multilinear polynomial $g(z)=\sum_{|S|\le3}a_Sz^S$
of norm at most $G$ on $[-1,1]^n$,
\[
 \sum_{1\le|S|\le3}|a_S|^{3/2}\le1024\,G^{3/2}.
 \label{eq:bh}\tag{CB}
\]
\end{lemma}
\begin{proof}
The case $G=0$ is immediate. Write $g=a_\emptyset+Q_1+Q_2+Q_3$
into homogeneous parts. Its odd part $O=Q_1+Q_3$ has norm at most
$G$, and
\[
 Q_3(z)=\tfrac43 O(z)-\tfrac83 O(z/2),\qquad Q_1=O-Q_3.
\]
Thus $\|Q_3\|\le4G$, $\|Q_1\|\le5G$. The even part gives
$\|Q_2\|\le2G$, because $|a_\emptyset|=|g(0)|\le G$.
In particular $\sum_{|S|=1}|a_S|\le5G$.

We supply the required array inequalities and moment estimates.
For a matrix $A$, let $R_i=(\sum_j A_{ij}^2)^{1/2}$ and
$S_j=(\sum_i A_{ij}^2)^{1/2}$. Three-factor Holder applied to
\[
 |A_{ij}|^{4/3}=
 (|A_{ij}|^2/R_i)^{1/3}(|A_{ij}|^2/S_j)^{1/3}(R_iS_j)^{1/3}
\]
gives
\[
 \left(\sum_{ij}|A_{ij}|^{4/3}\right)^{3/4}
 \le\left(\sum_iR_i\right)^{1/2}
      \left(\sum_jS_j\right)^{1/2}.\tag{M2}
\]
Zero slices can be removed. For a three-index array, put
$R_i=(\sum_{jk}A_{ijk}^2)^{1/2}$, with $S_j,T_k$ defined analogously.
Four-factor Holder, using the factorization
\[
 |A_{ijk}|^{3/2}=
 (|A_{ijk}|^2/R_i)^{1/4}
 (|A_{ijk}|^2/S_j)^{1/4}
 (|A_{ijk}|^2/T_k)^{1/4}(R_iS_jT_k)^{1/4},
\]
gives
\[
 \left(\sum_{ijk}|A_{ijk}|^{3/2}\right)^{2/3}
 \le\left(\sum_iR_i\sum_jS_j\sum_kT_k\right)^{1/3}.\tag{M3}
\]

For a linear form $Z$ in independent uniform signs,
$\E Z^4\le3(\E Z^2)^2$; expanding the fourth power proves this.
For a bilinear form $Z=\sum_{jk}b_{jk}\varepsilon_j\eta_k$ in two
independent sign groups,
\[
 \E Z^4\le9\left(\sum_{jk}b_{jk}^2\right)^2.
\]
Indeed, condition first on $\varepsilon$ and apply the linear
fourth-moment bound; expand the resulting square of a sum of squared
linear forms and apply Cauchy--Schwarz and the same bound once more.
Holder in the form
$\E Z^2\le(\E|Z|)^{2/3}(\E Z^4)^{1/3}$ therefore gives
$\|Z\|_2\le\sqrt3\|Z\|_1$ in the linear case and
$\|Z\|_2\le3\|Z\|_1$ in the bilinear case.

For $B(x,y)=\sum A_{ij}x_iy_j$, apply the linear estimate rowwise:
\[
 \sum_iR_i\le\sqrt3\E_y\sum_i|\sum_jA_{ij}y_j|
 \le\sqrt3\|B\|_\infty.
\]
The same holds for columns. Thus (M2) gives
$\|A\|_{4/3}\le\sqrt3\|B\|_\infty$.
For $T(x,y,z)=\sum A_{ijk}x_iy_jz_k$, the bilinear estimate likewise
gives $\sum_iR_i\le3\|T\|_\infty$ in each of the three directions;
(M3) gives $\|A\|_{3/2}\le3\|T\|_\infty$.

For $Q_2$, set $A_{ij}=a_{\{i,j\}}/2$ off the diagonal and zero
on it. Polarization gives
$B(x,y)=Q_2((x+y)/2)-Q_2((x-y)/2)$, hence $\|B\|\le4G$.
Therefore the quadratic coefficient vector has $\ell^{4/3}$ norm
at most $2^{1/4}4\sqrt3G$. Its $\ell^{3/2}$ norm is no larger, so
\[
 \sum_{|S|=2}|a_S|^{3/2}
 \le(2^{1/4}4\sqrt3G)^{3/2}<24G^{3/2}.
\]
The last strict inequality is equivalent to $2^3<3^2$.

For $Q_3$, set $A_{ijk}=a_{\{i,j,k\}}/6$ for distinct indices,
zero otherwise. The symmetric trilinear form has diagonal $Q_3$,
and polarization gives
\[
 T(x,y,z)=\frac1{2^3\,3!}
 \sum_{\varepsilon\in\{-1,1\}^3}
 \varepsilon_1\varepsilon_2\varepsilon_3
 Q_3(\varepsilon_1x+\varepsilon_2y+\varepsilon_3z).
\]
Homogeneity bounds its norm by $(3^3/3!)\|Q_3\|\le18G$.
Consequently
\[
 \sum_{|S|=3}|a_S|^{3/2}
 =\sqrt6\sum_{ijk}|A_{ijk}|^{3/2}
 \le\sqrt6(54G)^{3/2}=972G^{3/2}.
\]
The linear contribution is at most $(5G)^{3/2}<12G^{3/2}$.
Adding the three estimates gives less than $1008G^{3/2}$, which
implies the stated, slightly looser $1024G^{3/2}$ bound.
\end{proof}

\section{Extracting a physical influential coordinate}

Apply Lemmas~\ref{lem:generator}--\ref{lem:bh} to obtain
\[
 \sum_{1\le|S|\le3}|a_S|^{3/2}\le B.
\]
Let $M=\max_{1\le|S|\le3}|a_S|$. The variance assumption and
Lemma~\ref{lem:energy} imply
\[
 v\le L\sum_{S\ne\emptyset}a_S^2
 \le LB M^{1/2},\qquad
 M^2\ge\left(\frac{v}{LB}\right)^4.\tag{E}
\]
Choose a maximizing $S$, of size $k\in\{1,2,3\}$, and any $i\in S$.
Iterate \eqref{eq:lower} on $\partial_i p$ over the distinct
coordinates in $S\setminus\{i\}$. All expectations remain under $\mu$.
Then Jensen yields
\[
 \E(\partial_i p)^2
 \ge D^{-(k-1)}\sigma_{S\setminus i}^2\E(\partial_Sp)^2
 \ge D^{-(k-1)}\sigma_{S\setminus i}^2(\E\partial_Sp)^2.
\]
Since $\sigma_i^2\le1$, the same coordinate satisfies
\[
 \Inf_i^{\rm fl}\ge4D^{-(k-1)}a_S^2\ge4D^{-2}M^2,
 \qquad
 \Inf_i^{\rm pv}\ge D^{-k}a_S^2\ge D^{-3}M^2.
\]
For the second inequality also use
$\Inf_i^{\rm pv}\ge D^{-1}\sigma_i^2\E(\partial_i p)^2$.
Substituting (E) and $L=6C^3D^2$ proves
\[
 \Inf_i^{\rm fl}\ge\frac{D^{-10}C^{-12}}{324B^4}v^4,
 \qquad \Inf_i^{\rm pv}\ge\frac{D^{-11}C^{-12}}{1296B^4}v^4.
\]
Since $D=e^{4K}$, $C=(1-r)^{-1}$, monotonicity in $K,r$ proves
Theorem~\ref{thm:cubic}, including the resampling conversion.

Equation (DS) proves Corollary~\ref{cor:omega}. Pinning restricts
the set of boundary configurations, so its TV matrix is bounded
entrywise by the original matrix on the remaining coordinates.
Interaction supports shrink; after combining identical residual
terms by the triangle inequality, neither $K$ nor $\omega$ increases.
Residual singleton terms become unrestricted fields. Thus the
theorem itself is stable under pinnings, with the variance hypothesis
checked under the resulting law.

For Corollary~\ref{cor:pairwise}, the exact conditional logistic
formula gives $c_{ij}\le\tanh(\beta|J_{ij}|)$. If
$t_j=\tanh(\beta|J_{ij}|)$ and $\sum_jt_j<1$, the positive power
series of $\operatorname{arctanh}$ gives
\[
 \sum_j\beta|J_{ij}|=\sum_j\operatorname{arctanh}t_j
 \le\operatorname{arctanh}\Big(\sum_jt_j\Big)
 \le\tfrac12\log\frac{2-\kappa}{\kappa}.
\]
Substitute this bound for $K$ and $r\le1-\kappa$ in the theorem.

\section{A fixed-degree corollary and its precise limitation}

The cubic proof already supplies all its scalar coefficient estimates.
For completeness, there is a direct generalization using one imported
theorem. Defant--Mastylo--Perez, \emph{On the Fourier spectrum of
functions on Boolean cubes}, arXiv:1706.03670v1, equation (2) and
Theorem 1.1, state that for every integer $d\ge1$ there is a finite
$b_d$, independent of $n$, such that every real multilinear polynomial
$g$ of degree at most $d$ satisfies
\[
 \left(\sum_{|S|\le d}|\widehat g(S)|^{q_d}\right)^{1/q_d}
 \le b_d\|g\|_{\infty,\{-1,1\}^n},\qquad
 q_d=\frac{2d}{d+1}.
 \tag{BH}
\]
Here these are ordinary uniform Walsh coefficients of $g$, not
coefficients of $p$ in a presumed orthogonal Gibbs basis. This is the
only imported analytic theorem needed for the following corollary;
it is not needed for the cubic theorem.

\begin{corollary}[Every fixed degree]\label{cor:general}
Fix $d\ge1$, $v>0$, $K_0<\infty$ and $r_0<1$ before $n$. Under the
same Gibbs assumptions, write
\[
 C_0=(1-r_0)^{-1},\quad D_0=e^{4K_0},\quad
 L_d=d!C_0^dD_0^{d-1},\quad
 B_d=(b_dT_d(2))^{q_d},
\]
where $T_d$ is the Chebyshev polynomial. Every bounded real $p$ of
degree at most $d$ and variance at least $v$ has a coordinate satisfying
\[
 \Inf_i^{\rm fl}(p)\ge
 4D_0^{-(d-1)}\left(\frac{v}{L_d B_d}\right)^{d+1},\qquad
 \Inf_i^{\rm pv}(p)\ge
 D_0^{-d}\left(\frac{v}{L_d B_d}\right)^{d+1}.
\]
In particular the $n$-independent function requested in the prompt
exists uniformly for every fixed $\omega_0<1$.
\end{corollary}
\begin{proof}
Iterate \eqref{eq:upper} and
$\E f^2=(\E f)^2+\Var f$, using the affine bound at the last
step. For $k<d$ the resulting coefficient of
$\sum_{|S|=k}a_S^2$ is $k!C^kD^k$; at $k=d$ it is
$d!C^dD^{d-1}$. Each is at most $L_d$, including $d=1$.
Thus $\Var p\le L_d\sum_{S\ne\emptyset}a_S^2$.

The same generator has norm at most $T_d(2)$. Here is an elementary
verification: interpolate any degree-at-most-$d$ polynomial $h$
bounded by one on $[-1,1]$ at $t_j=\cos(j\pi/d)$, $0\le j\le d$.
The Lagrange weights at $2$ have signs $(-1)^j$. Since
$T_d(t_j)=(-1)^j$, their absolute sum is $T_d(2)$, and hence
$|h(2)|\le T_d(2)$. Apply this to $h(t)=p(ty/2)$ as before.
Applying (BH) to this real multilinear generator proves
$\sum_{S\ne\emptyset}|a_S|^{q_d}\le B_d$.

Since $2-q_d=2/(d+1)$, if $M=\max_{S\ne\emptyset}|a_S|$ then
$v\le L_dB_dM^{2/(d+1)}$, so
$M^2\ge(v/(L_dB_d))^{d+1}$. The identical derivative extraction
on a maximizing subset of size at most $d$ gives the two bounds.
\end{proof}

This corollary resolves the literal fixed-degree, dimension-free
target stated in the prompt and handoff. It must not be advertised
as a bound $(v/d)^{O(1)}$ with an absolute exponent. Its variance
exponent is $d+1$, and $L_d$ contains $d!$ while
$T_d(2)=((2+\sqrt3)^d+(2-\sqrt3)^d)/2$ grows exponentially.
The method therefore leaves the stronger joint-polynomial target
open. This quantitative loss is the remaining roadblock in this
method, after the cubic centered-Gram roadblock has been bypassed.

\section{Checks, scope, and what remains open}

For the pair-block cubic above, all $\sigma_i=1$, and its generator is
exactly
\[
 g_m(z)=\rho z_0+\frac1m\sum_{j=1}^m z_0z_{2j-1}z_{2j}.
\]
Its norm is $1+\rho$, and its expected-derivative energy is
$\rho^2+1/m$, rather than $1/m$. This explains explicitly why the
new energy estimate survives the old obstruction. It does not imply
the false base inequality (B) or uncentered marginal-product moment
domination from the earlier reports.

The hypotheses allow biased product laws ($K=r=0$), arbitrarily
large finite fields, and zero-field odd-arity interactions whose
one-site means need not vanish. Neither Boolean-valuedness nor
bounded support is used. Exactly deterministic product coordinates
can be removed before application. The general theorem is stated
for full-support Gibbs laws, and does not silently assert a Dobrushin
specification on boundaries of probability zero.

The companion \texttt{verify.py} uses Python's standard library and
the existing quadratic script's finite-cube helpers. Rational Gibbs
weights $\prod_A(1+t_A\chi_A)\prod_i(1+u_ix_i)$ represent exactly
$\beta J_A=\operatorname{arctanh}t_A$ and
$\beta h_i=\operatorname{arctanh}u_i$. Moments, derivatives, TV
coefficients, weighted squared coefficients, and the successive
Poincare/energy/extraction inequalities are checked rationally.
In particular $D=e^{4K}$ is computed as the exact rational maximum
$\max_i\prod_{A\ni i}((1+|t_A|)/(1-|t_A|))^2$.
Only fractional powers, displayed decimals, and the supplementary
weighted-generator evaluations use floating point. The script also
checks the generating-polynomial identity with rational shifts,
the interpolation formula, the exact pair-block cubic, signed
hyperedges, strong biases, pinnings, and small higher-degree cases
for the energy recurrence. Recorded output accompanies the source.
None of these finite tests proves a dimension-uniform statement.

The inverse-temperature condition is sufficient, not necessary:
if $W=\max_i\sum_{A\ni i}(|A|-1)|J_A|>0$, this proof applies
for $\beta W\le\omega_0<1$. If $W=0$ the law is product and all
finite $\beta$ are covered. No universal sharp $\beta$ threshold,
optimal influence exponent, all-temperature theorem, or new
efficient classical-simulation runtime is obtained here.
The supplied \texttt{gibbs\_history\_aa\_corrected.pdf}, Theorem 4.2
and Section 5, treats query simulation separately; its complete-history
independence must not be mistaken for conditional independence given
the final input. No such independence is used above.

\appendix
\section{Source and dependency ledger}
\begin{enumerate}
\item \path{prompt/polynomial-aa-next-agent-prompt.md};\\
\path{prompt/polynomial-aa-research-handoff.md}:\\
problem, conventions,
counterexample tests and proof standards. Their old status statements
are not assumed as mathematical inputs. Their referenced
\texttt{.ultramath/} sources are unavailable in this workspace.
\item \texttt{results/quadratic\_gibbs\_centered/report.tex}:
preceding quadratic argument, local estimates, and obstructions.
The analytic facts needed for this new cubic theorem are re-proved
above. Its \texttt{verify.py} supplies only finite enumeration helpers.
\item \texttt{gibbs\_history\_aa\_corrected.pdf}, particularly
Theorem 4.2 and Section 5: scope of the separate query result.
No theorem from it is imported into the influence proof.
\item A. Defant, M. Mastylo and A. Perez,
\emph{On the Fourier spectrum of functions on Boolean cubes},
\href{https://arxiv.org/pdf/1706.03670}{arXiv:1706.03670v1},
equation (2), p.~2, and Theorem 1.1, p.~3: scalar Boolean
Bohnenblust--Hille inequality (BH). Its exact real-valued,
degree-at-most-$d$, dimension-free statement was checked in the
primary paper. It is used only in Corollary~\ref{cor:general}.
The cubic coefficient lemma is proved independently above.
\end{enumerate}
\end{document}
